% Generated draft: requires mathematical and source review.
\documentclass[11pt]{article}
\usepackage[margin=25mm]{geometry}
\usepackage{fontspec}
\setmainfont{Latin Modern Roman}
\usepackage{amsmath,amssymb,mathtools,unicode-math}
\setmathfont{Latin Modern Math}
\usepackage{enumitem,needspace,xurl,hyperref}
\setlength{\parindent}{0pt}
\setlength{\parskip}{5pt}
\setlength{\emergencystretch}{3em}
\newcommand{\N}{\mathbb{N}}
\newcommand{\Z}{\mathbb{Z}}
\newcommand{\Q}{\mathbb{Q}}
\newcommand{\R}{\mathbb{R}}
\newcommand{\C}{\mathbb{C}}
\newcommand{\F}{\mathbb{F}}
\newcommand{\A}{\mathbb{A}}
\newcommand{\PP}{\mathbb{P}}
\newcommand{\E}{\mathbb{E}}
\newcommand{\eps}{\varepsilon}
\newcommand{\dd}{\,\mathrm{d}}
\begin{document}
\section*{UnsolvedMath ID: 30001715}
\textbf{Source problem code:} catalogue entry
\subsection{Definitions and mathematical statement}
Write a quiver as \(\Gamma=(I,H,s,t)\), where \(I\) is its vertex set, \(H\) its arrow set, and \(s,t:H\to I\) give the source and target of each arrow. Write \(v=(v_i)_{i\in I}\) for the dimension vector appearing in the source. The source does not specify finiteness assumptions, restrictions on loops or multiple arrows, or whether zero coordinates are permitted.

The supplied record proposes the following program for every quiver \(\Gamma\) and dimension vector \(v\): construct a projective curve \(C_{\Gamma,v}\) with one rational component for each vertex \(i\), having multiplicity \(v_i\), and with nodal intersections prescribed by the edges; form a compactified Jacobian \(J_v\) of this curve; and seek a cohomological interpretation of \(A_\Gamma(v,q)\) through \(J_v\) and an associated character variety.

This is an incompletely specified research program. The base field, the scheme structure realizing the component multiplicities, and the precise incidence construction are not supplied. In particular, for multiplicities greater than one, it is unclear whether the nodal condition concerns the reduced support or the full curve. The treatment of loops, multiple arrows, and zero coordinates is unspecified. The compactified Jacobian requires a precise moduli definition and any necessary degree, stability, and polarization choices. Neither the associated character variety nor its parameters are defined. The title suggests that \(A_\Gamma(v,q)\) denotes a Kac polynomial, but its definition and the domain of \(q\) are not supplied. No cohomology theory, grading, normalization, or proposed identity is given. Consequently, the record does not determine a precise universally quantified theorem or conjectural equality.
\subsection{Short English statement}
For each quiver and dimension vector, the source proposes relating a curve assembled from rational components, its compactified Jacobian, an associated character variety, and the indicated Kac polynomial. Essential constructions and the intended cohomological relation remain unspecified, so this is a statement-only draft of a research program requiring mathematical review.
\subsection{Sources}
\begin{itemize}
\item Supplied local UnsolvedMath registry record: permanent ID \(30001715\), source code \texttt{catalogue entry}, title ``Compactified Jacobians and Kac Polynomials of Quivers,'' in the Oberwolfach Reports Open Problems collection. The recorded proposed year is 2011; no individual proposer is identified.
\item \emph{Representation Theory of Quivers and Finite Dimensional Algebras} (2011), Oberwolfach Reports, pp.\ 523--608, DOI \texttt{10.4171/owr/2011/10}: \url{https://doi.org/10.4171/owr/2011/10}. This is also the supplied literature source labeled ``Oberwolfach Report 10/2011.'' No specific page, contribution, or passage within the report is identified.
\item Literature assessment dated exactly \texttt{2026-08-21}: the supplied curation labels the problem open, describes it as contextual material for a proposed compactified-Jacobian and character-variety program rather than a fully stated theorem, and reports that no complete general cohomological interpretation for the extracted arbitrary-quiver setting was located. This is supplied dated curation, not an independently verified literature assessment; its time-sensitive status and cited claims require independent verification.
\item The registry labels the statement \texttt{corrected\_verified} and reports a manual semantic, typographic, and LaTeX review claiming that the cleanup is self-contained and introduces no unsupported source-dependent content. That supplied editorial label does not resolve the formulation gaps above or establish mathematical correctness.
\end{itemize}
\end{document}
